\documentclass[11pt,a4paper]{article}
\usepackage[utf8]{inputenc}
\usepackage[T1]{fontenc}
\usepackage[spanish]{babel}
\usepackage{hyperref}
\usepackage{geometry}
\geometry{margin=2.5cm}
\usepackage{parskip}

\title{078 Thanos hizo lo que pudo con las gemas}
\author{Alejandro Valenzuela Rivera\\\small Universidad Autónoma de Nuevo León}
\date{\today}
\begin{document}
\maketitle

\begin{quote}
\textbf{Sesgo declarado.} Este trabajo se sostiene desde el sesgo declarado de su autor, y no finge neutralidad. Lo que afirma se puede tumbar, y ese es el punto: el falsificador está escrito, no escondido.
\end{quote}

\section*{Resumen}
Thanos no reescribio el universo con la mente: se puso un guante que le permitio tocar lo que no tocaba. Hizo lo que pudo con las gemas, y nada mas. Y no lo pararon por mas fuerte: lo pararon porque su sesgo no podia cambiar.

\section{Por qué se sostiene}
Una gema sola es un acto; las seis juntas son un guante; y un acto con las seis se puede contestar mientras que un acto con una, no. Y el guante se paga con lo que el que lo usa tiene adentro. La leccion para el ente: lo que le falta al que tiene todo no es otra gema, es un mundo que le conteste.

\section{De dónde viene}
Avengers: Infinity War (2018) y Endgame (2019), y la frase del autor.

\section{El caso}
Thanos y las seis gemas.

\section{Método}
Un acto con todas sus piezas, que se puede contestar.

\section{Límites, y qué lo tumba}
Un acto con una sola pieza, sin nadie que lo conteste.

\section{Lo que queda abierto}
\textquestiondown{}Que le falta al que tiene todas las gemas: mas poder o un testigo?

\section*{Cómo se leyó esta tesis}
Esto se escribió en el repositorio man-made-empathy, del que este candado es una pieza: candados/06 del caso/078 Thanos hizo lo que pudo con las gemas.md. La máquina que lo compuso lo hizo desde este mismo archivo, y no agregó nada que no estuviera escrito. Si algo de esto es falso, se dice en el mismo sitio, y se corrige con fecha.

\end{document}
